\begin{definition}[\texttt{eVariationOn}, Mathlib]
  For $f : \mathbb{R} \to E$ into a pseudo-emetric space and $s \subseteq \mathbb{R}$,
  \texttt{eVariationOn f s} is the extended-real total variation of $f$ on $s$:
  the supremum, over finite increasing chains $x_0 \le x_1 \le \dots \le x_n$ in $s$, of
  $\sum_i \mathrm{edist}(f(x_i), f(x_{i+1}))$. This is Mathlib's formalization of the
  length of the (possibly non-injective) path traced by $f$ restricted to $s$.
\end{definition}

\begin{definition}[\texttt{HasConstantSpeedOnWith} / \texttt{HasUnitSpeedOn}, Mathlib]
  \texttt{HasConstantSpeedOnWith f s l} means: for all $x,y \in s$,
  $\texttt{eVariationOn}\, f\, (s \cap [x,y]) = l \cdot (y-x)$ (as extended reals, via
  \texttt{ENNReal.ofReal}). \texttt{HasUnitSpeedOn f s} is the case $l=1$, i.e.\ $f$
  traces out length exactly $y-x$ on every subinterval $[x,y]$ of $s$ -- Mathlib's
  formalization of ``$f$ is parametrized by arc length on $s$.''
\end{definition}

\begin{proposition}[\texttt{eVariationOn.edist\_le}, Mathlib]
  If $x, y \in s$ then $\mathrm{edist}(f(x),f(y)) \le \texttt{eVariationOn}\, f\, (s \cap [x,y])$
  (with $[x,y]$ meaning $[\min(x,y),\max(x,y)]$ as needed): the displacement between two points
  of $f$'s image is bounded by the variation of $f$ between them.
\end{proposition}

\begin{definition}[\texttt{LipschitzWith}, Mathlib]
  \texttt{LipschitzWith K f} for $K : \mathbb{R}_{\ge 0}$ means
  $\mathrm{edist}(f(x),f(y)) \le K \cdot \mathrm{edist}(x,y)$ for all $x,y$ in the domain of $f$.
\end{definition}

\begin{definition}[\texttt{Measure.map}, Mathlib]
  For a (possibly non-measurable) function $f : X \to Y$ between measurable spaces and a measure
  $\mu$ on $X$, \texttt{Measure.map f} $\mu$ is the pushforward measure: when $f$ is measurable,
  $(\texttt{Measure.map}\, f\, \mu)(B) = \mu(f^{-1}(B))$ for Borel $B \subseteq Y$; when $f$ is not
  measurable the pushforward is defined to be the zero measure (a junk value that is never invoked
  in this tablet, since every curve's underlying map is shown Lipschitz, hence measurable).
\end{definition}

\begin{definition}[\texttt{deriv}, Mathlib]
  For $f : \mathbb{R} \to \mathbb{R}$ and $x \in \mathbb{R}$, \texttt{deriv f x} is the derivative
  of $f$ at $x$ if $f$ is differentiable at $x$, and the junk value $0$ otherwise.
\end{definition}

\begin{definition}[interval integral, Mathlib]
  For $f : \mathbb{R} \to \mathbb{R}$ and $a, b \in \mathbb{R}$, the notation
  $\int_a^b f(x)\,dx$ (Mathlib: \texttt{intervalIntegral}, written \texttt{$\int$ x in a..b, f x})
  denotes the signed Lebesgue integral of $f$ on $[a,b]$: it equals
  $\int_{[a,b]} f\,d\mathrm{Leb}$ when $a \le b$, and $-\int_{[b,a]} f\,d\mathrm{Leb}$ when $b < a$;
  it is $0$ (a junk value) when $f$ is not integrable on the relevant interval.
\end{definition}

\begin{definition}[\texttt{AbsolutelyContinuousOnInterval}, Mathlib]
  For $f : \mathbb{R} \to X$ into a (pseudo)metric space and $a,b \in \mathbb{R}$,
  \texttt{AbsolutelyContinuousOnInterval f a b} means: $f$ is absolutely continuous on
  $[\min(a,b),\max(a,b)]$, i.e.\ for every $\varepsilon>0$ there is $\delta>0$ such that for every
  finite collection of pairwise-disjoint subintervals $(a_i,b_i)$, $i<n$, of $[\min(a,b),\max(a,b)]$
  with $\sum_i d(a_i,b_i) < \delta$, we have $\sum_i d(f(a_i),f(b_i)) < \varepsilon$.
\end{definition}

\begin{proposition}[\texttt{LipschitzOnWith.absolutelyContinuousOnInterval}, Mathlib]
  If $f$ is $K$-Lipschitz on $[\min(a,b),\max(a,b)]$ then
  $\texttt{AbsolutelyContinuousOnInterval}\ f\ a\ b$.
\end{proposition}

\begin{proposition}[\texttt{AbsolutelyContinuousOnInterval.ae\_differentiableAt}, Mathlib]
  If $f : \mathbb{R} \to \mathbb{R}$ satisfies $\texttt{AbsolutelyContinuousOnInterval}\ f\ a\ b$,
  then $f$ is differentiable at $x$ for a.e.\ $x \in [\min(a,b),\max(a,b)]$.
\end{proposition}

\begin{proposition}[\texttt{AbsolutelyContinuousOnInterval.intervalIntegrable\_deriv}, Mathlib]
  If $f : \mathbb{R} \to \mathbb{R}$ satisfies $\texttt{AbsolutelyContinuousOnInterval}\ f\ a\ b$,
  then $x \mapsto \texttt{deriv}\ f\ x$ is interval-integrable on $a..b$.
\end{proposition}

\begin{proposition}[\texttt{HasDerivAt.le\_of\_lipschitz}, Mathlib]
  If $f : \mathbb{R} \to \mathbb{R}$ has derivative $f'(x_0)$ at $x_0$ (in the sense of
  \texttt{HasDerivAt}) and $f$ is $K$-Lipschitz (on all of $\mathbb{R}$), then $|f'(x_0)| \le K$.
  In particular, if $f$ is differentiable at $x_0$ and $K$-Lipschitz, $|\texttt{deriv}\ f\ x_0| \le K$.
\end{proposition}

\begin{proposition}[\texttt{intervalIntegral.integral\_add\_adjacent\_intervals}, Mathlib]
  If $f : \mathbb{R} \to \mathbb{R}$ is interval-integrable on $a..b$ and on $b..c$, then
  $\int_a^b f + \int_b^c f = \int_a^c f$.
\end{proposition}

\begin{proposition}[\texttt{intervalIntegral.integral\_comp\_sub\_left}, Mathlib]
  For $f : \mathbb{R} \to \mathbb{R}$ and $a,b,d \in \mathbb{R}$,
  $\int_a^b f(d - x)\,dx = \int_{d-b}^{d-a} f(x)\,dx$.
\end{proposition}

\begin{proposition}[\texttt{intervalIntegral.norm\_integral\_le\_of\_norm\_le\_const}, Mathlib]
  For $f : \mathbb{R} \to \mathbb{R}$ and $a \le b$, if $|f(x)| \le C$ for all $x \in [a,b]$, then
  $\left|\int_a^b f(x)\,dx\right| \le C \cdot (b-a)$.
\end{proposition}

\begin{proposition}[\texttt{AbsolutelyContinuousOnInterval.integral\_deriv\_eq\_sub}, Mathlib]
  If $f : \mathbb{R} \to \mathbb{R}$ satisfies $\texttt{AbsolutelyContinuousOnInterval}\ f\ a\ b$,
  then $\int_a^b \mathrm{deriv}(f)(x)\,dx = f(b) - f(a)$ (the Fundamental Theorem of Calculus for
  absolutely continuous functions).
\end{proposition}

\begin{proposition}[\texttt{deriv\_comp\_const\_sub}, Mathlib]
  For $f : \mathbb{R} \to \mathbb{R}$ and $a,x \in \mathbb{R}$,
  $\mathrm{deriv}(t \mapsto f(a-t))(x) = -\mathrm{deriv}(f)(a-x)$, unconditionally (including when $f$ fails to be
  differentiable at $a-x$, in which case both sides are $0$ by the junk-value convention).
\end{proposition}

\begin{proposition}[\texttt{integral\_id}, Mathlib]
  For $a,b \in \mathbb{R}$, $\int_a^b x\,dx = (b^2 - a^2)/2$; in particular, at $a=0$,
  $\int_0^b x\,dx = b^2/2$.
\end{proposition}

\begin{proposition}[\texttt{measurable\_of\_tendsto\_metrizable}, Mathlib]
  Let $\beta$ be a (pseudo-)metrizable Borel space and $(f_n)_{n\in\mathbb{N}}$, $g$ functions
  $\alpha \to \beta$. If each $f_n$ is measurable and $f_n \to g$ pointwise (equivalently,
  $\texttt{Filter.Tendsto}\ f\ \texttt{Filter.atTop}\ (\mathcal{N}\ g)$ in the topology of
  pointwise convergence on $\alpha \to \beta$), then $g$ is measurable.
\end{proposition}

\begin{definition}[\texttt{Measure.infinitePi}, Mathlib]
  For an index type $\iota$ and a family of probability measures $\mu_i$ on measurable spaces
  $X_i$, $\iota \in \iota$, \texttt{Measure.infinitePi} $\mu$ is the unique probability measure
  $P$ on $\prod_{i} X_i$ such that, for every finite $s \subseteq \iota$ and measurable
  rectangle $\prod_{i \in s} t_i \times \prod_{i \notin s} X_i$ (with $t_i$ measurable),
  $P$ assigns it mass $\prod_{i \in s} \mu_i(t_i)$: the countable (or arbitrary) product measure.
\end{definition}

\begin{proposition}[\texttt{measurePreserving\_eval\_infinitePi}, Mathlib]
  For each coordinate $i$, the evaluation map $\mathrm{ev}_i \colon \prod_j X_j \to X_i$,
  $\mathrm{ev}_i(\omega) = \omega_i$, is measure-preserving from $\texttt{Measure.infinitePi}\ \mu$
  to $\mu_i$.
\end{proposition}

\begin{proposition}[\texttt{ProbabilityTheory.iIndepFun\_infinitePi}, Mathlib]
  For any family of measurable maps $X_i \colon \Omega_i \to \beta_i$ (in particular the identity
  maps used here), the coordinate maps $\omega \mapsto X_i(\omega_i)$, $i \in \iota$, are jointly
  independent under $\texttt{Measure.infinitePi}\ \mu$.
\end{proposition}

\begin{proposition}[\texttt{ProbabilityTheory.strong\_law\_ae\_real}, Mathlib]
  Let $(Y_i)_{i\in\mathbb{N}}$ be a pairwise-independent, identically-distributed sequence of
  real-valued integrable random variables on a probability space $(\Omega,P)$. Then, for
  $P$-almost every $\omega$, $\frac{1}{n}\sum_{i=0}^{n-1} Y_i(\omega) \to \mathbb{E}_P[Y_0]$ as
  $n \to \infty$.
\end{proposition}

\begin{proposition}[\texttt{MeasureTheory.ae\_all\_iff}, Mathlib]
  For a countable index type $\iota$ and a family of predicates $p_i \colon \alpha \to
  \texttt{Prop}$, $(\forall^{a.e.} x,\ \forall i,\ p_i(x)) \iff (\forall i,\ \forall^{a.e.} x,\
  p_i(x))$: in particular, a countable intersection of full-measure sets is full-measure.
\end{proposition}
